\section{Metodologiyalar va Qiyosiy Matritsa}
Tahlil qilingan ilmiy manbalarning umumiy qiyosiy matritsasi 1-jadvalda keltirilgan.

\begin{table}[htbp]
\centering
\caption{AI Agents for Scientific Discovery bo'yicha metodologiyalar solishtirma tahlili}
\begin{tabular}{llll}
\toprule
\textbf{Mualliflar va Yil} & \textbf{Yondashuv} & \textbf{Ma'lumotlar To'plami} & \textbf{Asosiy Natija} \\
\midrule
Vaswani et al. (2017) & Transformer & WMT 2014 & Self-Attention asosida SOTA \\
Devlin et al. (2019) & BERT & BooksCorpus, Wiki & Ikki tomonlama kontekst \\
Tadqiqotchi et al. (2024) & Hybrid Architecture & Domain Benchmark & Resurs tejamkorligi ($+18\%$) \\
\bottomrule
\end{tabular}
\end{table}
